\subsection{Initial-visit MC-O-PI converges with uniform updates}
This section proves that stochastic perturbations in initial-visit MC-O-PI cannot consistently offset the policy improvement of mean-field solutions. As a result, the algorithm converges to $\qopt$ when updates are uniform over the actions of each state.

\begin{theorem}
\label{thm: convergence mcopi uniform}
Under \autoref{asp: standing}, the action-value targets $(q_{\pi_k})$ change only finitely many times almost surely, and
\[
q_k\longrightarrow q_{\pi_*}\qquad\text{almost surely}.
\]
Every sufficiently late selected policy is optimal. If every state has a unique optimal action, the selected policy is eventually constant.
\end{theorem}

\subsubsection{Proof strategy}
\label{sec: limitations of existing methods}
The mean-field argument improves the policy at every transition. A sampled update can instead worsen it, so we seek a fixed positive probability of improvement at each sufficiently late transition. Since only finitely many deterministic policies exist, repeated opportunities for improvement will then imply convergence.

This follows the combined stability--ODE strategy of \citet[Section~5.4]{Kushner2003StochasticApplications}. The difficulty is that a trajectory can enter a policy region $Q(\pi)$ arbitrarily close to its boundary. Recurrence to that region alone does not give repeated visits to the fixed interior subset used in a standard lock-in estimate; see the setup in \citet[Section~3]{Borkar2002OnApproximation}.
We therefore first create a positive separation between the currently selected action and actions that would worsen the policy, and then control the probability that noise removes that separation. A sufficiently large fixed multiple of the current learning rate is enough: the separation may tend to zero at successive entries. Condition~\eqref{eq:comparison} makes the resulting probability bound uniform over sufficiently late entries. It is a sufficient condition for this proof, not a claimed necessity for convergence of the algorithm.

\subsubsection{Actions that would worsen the policy}
Consider a policy $\pi$ selected at a time $t$. The pairs
\[
\mathcal B_\pi=\big\{(s,a)\in\mathcal S\times\mathcal A:
q_\pi(s,a)<q_\pi(s,\pi(s))\big\}
\]
would worsen its selected action according to $q_\pi$. If none of these pairs becomes greedy, then at the next transition the new policy satisfies the Policy Improvement Theorem. To express this condition in terms of the estimates, we compare their estimated action values.

\begin{definition}[Gap variables]
\label{def: gap variables}
Fix a deterministic bound $K>\max_{\pi\in\mathcal P}\|q_\pi\|$ and a deterministic restart $r$. Let
\[
\zeta_r=\inf\{k\ge r:\|q_k\|>K\}
\]
be the first time the estimate leaves this bound.
For a stopping time $t\ge r$, on $\{t<\infty,t<\zeta_r\}$ set $\pi=\pi_t$ and let
\[
\tau=\inf\{k>t:q_{\pi_k}\ne q_\pi\}
\]
be the next change of the action-value target.
For each pair $(s,a)\in\mathcal B_\pi$ and each $k\ge t$, define
\begin{align}
\label{eq:gaps}
d_k(s,a)&=q_k(s,\pi(s))-q_k(s,a),\\
d_\pi(s,a)&=q_\pi(s,\pi(s))-q_\pi(s,a).
\end{align}
The gap noise is $\xi_k(s,a)=w_k(s,\pi(s))-w_k(s,a)$.
\end{definition}

The target remains $q_\pi$ before $\tau$, even if the selected policy changes. Equal update probabilities within each state therefore give the gap recursion
\begin{equation}
\label{eq:gapdrift}
d_{k+1}(s,a)=d_k(s,a)+\alpha_k\big(\sigma_k(s)(d_\pi(s,a)-d_k(s,a))+\xi_k(s,a)\big).
\end{equation}
Thus, on average, the estimate gap is pulled towards the positive target gap $d_\pi(s,a)$. The noise has the following bounds for $t\le k<\tau\wedge\zeta_r$:
\begin{equation}
\label{eq:gapvariance}
\mathbb E[\xi_k(s,a)\mid\mathcal F'_k]=0,
\qquad
\mathbb E[\xi_k(s,a)^2\mid\mathcal F'_k]\le c_\xi\sigma_k(s),
\qquad c_\xi=2(c_v+4K^2).
\end{equation}
Indeed, only the initial pair is updated. The two possible updates forming this gap are disjoint, each has probability $\sigma_k(s)$, and before $\zeta_r$ the difference between an estimate and its target has magnitude at most $2K$. Their total second moment bounds the variance in \eqref{eq:gapvariance}.
The stopping time $\zeta_r$ is used only to justify these deterministic bounds. Boundedness will remove it at the end of the proof; we never condition the noise on future containment.

We first determine how large a positive gap must be for stochastic updates to preserve it with high probability. We then construct that margin through favorable updates, a step called \emph{seeding}, and combine the two estimates to obtain a positive probability of policy improvement.

Choose
\begin{equation}
\label{eq:constants}
\delta=\frac14\min_{\substack{\pi\in\mathcal P\\(s,a)\in\mathcal B_\pi}}d_\pi(s,a)>0.
\end{equation}
If the set in this minimum is empty, the Policy Improvement Theorem applies between every two policies in both directions. All policies then have the same optimal values, and \autoref{thm: suboptimal blocks finite} proves the theorem directly. We henceforth consider the nonempty case.

\subsubsection{Preserving positive gaps}
The following estimate bounds the probability that a positive gap closes before the next target change or exit. The drift and noise variance in \eqref{eq:gapdrift}--\eqref{eq:gapvariance} both carry the per-action update probability $\sigma_k(s)$ at state $s$. Measuring accumulated learning at that state will let us control this probability even when its updates are arbitrarily far apart. Write
\[
\bar\alpha_k=\sup_{j\ge k}\alpha_j
\]
for the largest remaining learning rate.

\begin{lemma}[Gap preservation]
\label{lem:barrier}
Consider a pair $(s,a)\in\mathcal B_\pi$ in \autoref{def: gap variables}, and a stopping time $t'\ge t$ with $t'<\infty$, $t'\le\tau\wedge\zeta_r$ and $d_{t'}(s,a)>0$.
Assume deterministic learning rates $0<\alpha_k\le1$ with $\sum_k\alpha_k^2<\infty$. Suppose that, under a given probability law, the gap recursion \eqref{eq:gapdrift} and conditional moments \eqref{eq:gapvariance} hold for $t'\le k<\tau\wedge\zeta_r$. Under that law,
\begin{equation}
\label{eq:barrier}
\mathbb P\!\left(\exists\text{ finite }n\in[t',\tau\wedge\zeta_r]:d_n(s,a)\le0\ \middle|\ \mathcal F'_{t'}\right)
\le \frac{4c_\xi}{3\delta}\frac{\bar\alpha_{t'}}{d_{t'}(s,a)}
+\frac{4c_\xi}{\delta^2}\sum_{k=t'}^\infty\alpha_k^2.
\end{equation}
The crossing event includes the update ending at $\tau\wedge\zeta_r$.
\end{lemma}
Equation~\eqref{eq:barrier} bounds the chance that one specified gap closes. If its starting value is bounded below by a fixed positive number, both terms become arbitrarily small as $t^{\prime}$ increases. If the starting gap also tends to zero, however, the first term depends on the ratio of the largest remaining learning rate to that gap: waiting alone need not make it small. A sufficiently large multiple of that learning rate controls this ratio, while a sufficiently late start makes the squared-learning-rate tail small. After proving the estimate, we construct such a margin through favorable updates.

\begin{proof}
Fix $(s,a)$ and write $d_k=d_k(s,a)$ within this proof. Stop the noise sum at the first gap closure or at the endpoint, whichever comes first:
\[
\eta=(\tau\wedge\zeta_r)\wedge\inf\{n\ge t':d_n\le0\},\qquad
M_n=\sum_{k=t'}^{n-1}\mathbf1_{\{k<\eta\}}\alpha_k\xi_k(s,a).
\]
The decision to retain increment $k$ is known from $\mathcal F'_k$, so $M_n$ has zero conditional-mean increments. Define the accumulated learning at state $s$ by
\begin{equation}\label{eq:clock}
A_n=\sum_{k=t'}^{n-1}\alpha_k\sigma_k(s).
\end{equation}
Its increments bound the conditional variance of the noise:
\[
\mathbb E[(M_{k+1}-M_k)^2\mid\mathcal F'_k]
\le c_\xi\alpha_k^2\sigma_k(s)
\le c_\xi\bar\alpha_{t'}(A_{k+1}-A_k).
\]
Suppose the first gap closure occurs at a finite $n\le\tau\wedge\zeta_r$. We distinguish whether the gap reached $\delta$ before closing.

\emph{A crossing before reaching $\delta$.} For $t'\le k<n$, the target gap is at least $4\delta$ while $0<d_k<\delta$. Summing \eqref{eq:gapdrift} therefore gives
\[
0\ge d_n\ge d_{t'}+3\delta A_n+M_n.
\]
Thus the noise must cross the growing threshold $b_n=d_{t'}+3\delta A_n$. To bound this event, divide each noise increment by the threshold for the next step:
\[
Z_n=\sum_{k=t'}^{n-1}\frac{M_{k+1}-M_k}{b_{k+1}}.
\]
The denominator $b_{k+1}$ is known from $\mathcal F'_k$, since $\sigma_k(s)$ is chosen before drawing the episode. Hence the rescaled increments still have conditional mean zero. To ensure square integrability when the initial gap is random, first restrict to $\{d_{t'}\ge\varepsilon\}$ for a fixed $\varepsilon>0$, an event known at $t'$. Equivalently, multiply the increments by its indicator. All the following conditional estimates hold on this event; letting $\varepsilon\downarrow0$ covers every positive initial gap. Random starts are treated on their finite-time events, as in the appendix.

Summation by parts gives
\[
M_n=b_nZ_n-\sum_{k=t'+1}^{n-1}(b_{k+1}-b_k)Z_k,
\qquad
|M_n|\le2b_n\max_{t'\le m\le n}|Z_m|,
\]
because the nonnegative differences $b_{k+1}-b_k$ sum to at most $b_n$. A crossing therefore requires $\sup_{n\ge t'}|Z_n|\ge1/2$. The conditional martingale maximal inequality recalled in the appendix bounds its probability by
\begin{align*}
4\,\mathbb E\!\left[
\sum_{k\ge t'}\frac{c_\xi\bar\alpha_{t'}(A_{k+1}-A_k)}{(d_{t'}+3\delta A_{k+1})^2}
\ \middle|\ \mathcal F'_{t'}\right]
&\le4c_\xi\bar\alpha_{t'}\int_0^\infty\frac{\mathrm dx}{(d_{t'}+3\delta x)^2}\\
&=\frac{4c_\xi\bar\alpha_{t'}}{3\delta d_{t'}}.
\end{align*}
The integral bound uses the fact that the integrand is decreasing: each summand without $c_\xi\bar\alpha_{t'}$ is at most its integral from $A_k$ to $A_{k+1}$. No divergence of $A_n$ is required.

\emph{A crossing after reaching $\delta$.} Let $\ell<n$ be the last index with $d_\ell\ge\delta$. Since $0\le\alpha_\ell\sigma_\ell(s)\le1$, the noise-free part of the update at $\ell$ is a convex combination of $d_\ell\ge\delta$ and $d_\pi(s,a)\ge4\delta$, so it is at least $\delta$. All subsequent gaps before $n$ lie in $(0,\delta)$ and have nonnegative mean increments. Consequently,
\[
0\ge d_n\ge\delta+M_n-M_\ell.
\]
This forces $\sup_{n\ge t'}|M_n|\ge\delta/2$. Since $\sigma_k(s)\le1$, the same maximal inequality bounds its conditional probability by
\[
\frac{4c_\xi}{\delta^2}\sum_{k\ge t'}\alpha_k^2.
\]
The last-visit index $\ell$ is used only in a pathwise inequality; we never condition at it. Both cases retain the increment ending at $\tau\wedge\zeta_r$. Adding their bounds proves \eqref{eq:barrier}.
\end{proof}

\subsubsection{Seed: creating positive gaps}
The preservation bound \eqref{eq:barrier} now gives us a concrete objective: lower-bound its denominator $d_{t'}(s,a)$ relative to the largest remaining learning rate. At a newly selected policy, some estimated gaps may be zero. We create the required margin through a finite number of favorable updates. States may wait arbitrarily long between updates, so each state completes this construction separately; its gaps will then be protected while other states are still waiting.

Suppose the reference action $\pi(s)$ is still selected and $d_k(s,a)<\delta$. Since
\[
\big(q_\pi(s,\pi(s))-q_k(s,\pi(s))\big)
+\big(q_k(s,a)-q_\pi(s,a)\big)
=d_\pi(s,a)-d_k(s,a)>3\delta,
\]
at least one of these two estimation errors exceeds $3\delta/2$. We request a favorable update as follows:
\begin{enumerate}
\item If $q_\pi(s,\pi(s))-q_k(s,\pi(s))>3\delta/2$, update $(s,\pi(s))$ and require $v_k(s,\pi(s))\ge-\delta/2$.
\item Otherwise update $(s,a)$ and require $v_k(s,a)\le\delta/2$; in this case $q_k(s,a)-q_\pi(s,a)>3\delta/2$.
\end{enumerate}
Either update increases the intended gap by more than $\delta\alpha_k$, leaves every other gap against the reference action unchanged or larger, and keeps the reference action greedy. The selection rule \eqref{eq:greedy-pol-tie-breaking} therefore retains it.

The conditional return moments ensure that each one-sided noise requirement has probability at least $\delta^2/(4c_v+\delta^2)$, by the one-sided Chebyshev inequality (the short conditional calculation is recalled in the appendix). Conditional on selecting the state, each action is sampled uniformly. Thus every favorable update has conditional probability at least
\begin{equation}
\label{eq:localprob}
p_0=\frac{\delta^2}{\max_{s\in\mathcal S}|\mathcal A(s)|\,(4c_v+\delta^2)}>0.
\end{equation}
No lower bound on the probability of selecting the state is needed.

\begin{definition}[Seed construction]
\label{def: seed construction}
For the gaps and stopping times in \autoref{def: gap variables}, fix a positive integer $N$. To give every gap an initial margin, order the actions in $\mathcal B_\pi$ at each state. Skip a pair if its gap has already reached $\delta$; otherwise request up to $N$ favorable updates for it, stopping earlier if it reaches $\delta$. Whenever an unfinished state is sampled, request the next favorable update on its list. Completed states receive ordinary updates. Stop the construction at the first failed request or at $\tau\wedge\zeta_r$.
Write $T_s$ for the iteration when state $s$ completes its list, with $T_s=t$ if it is complete initially and $T_s=\infty$ if the construction stops first. Skips are processed at $t$ and after every successful update, before the next sample, so $T_s$ is a stopping time.
\end{definition}

\begin{lemma}[Seed]
\label{lem:seeds}
Under \autoref{asp: standing}, consider the gaps and stopping times of \autoref{def: gap variables}, with $t\ge k_\alpha$, and let $N$ be a positive integer. The seed construction of \autoref{def: seed construction} requests at most $|\mathcal S\times\mathcal A|N$ favorable updates. Along successful requests, each unfinished state retains its reference action. If $T_s<\infty$, then
\begin{equation}
\label{eq:seedmargin}
d_{T_s}(s,a)\ge\min\{\delta,N\delta\bar\alpha_{T_s}/c_\alpha\}
\qquad\text{for every }a\text{ with }(s,a)\in\mathcal B_\pi.
\end{equation}
This includes a state completed at the endpoint $\tau\wedge\zeta_r$. Each request has conditional probability at least $p_0$ given the current history and the selection of its unfinished state.
\end{lemma}
\begin{proof}
A favorable update only raises the reference estimate or lowers a competing estimate. It keeps the reference action greedy, so \eqref{eq:greedy-pol-tie-breaking} retains that action. At a state not updated, the policy is unchanged. Induction therefore retains the reference action at every unfinished state, including during waits and after its last favorable update.

Every processed gap either reaches $\delta$ or receives $N$ increases of more than $\delta\alpha_k$. Later favorable updates at the same state cannot lower it. For each seeding step $k<T_s$, condition~\eqref{eq:comparison} gives $\alpha_j\le c_\alpha\alpha_k$ for every $j\ge T_s$, hence $\bar\alpha_{T_s}\le c_\alpha\alpha_k$. Summing these $N$ increases proves \eqref{eq:seedmargin}.
At most $|\mathcal S\times\mathcal A|$ pairs are processed, and the probability of each request is bounded below by \eqref{eq:localprob}.
\end{proof}
